\documentclass[11pt]{article}
\usepackage[margin=0.9in]{geometry}
\usepackage{amsmath,amssymb,amsthm}
\usepackage[colorlinks=true,urlcolor=blue,linkcolor=blue]{hyperref}
\newtheorem{theorem}{Theorem}
\title{Unimodality for one occurrence of the pattern 112\\\large MathDB 345816, Task 1}
\author{Principia Math In Progress}
\date{22 September 2026}
\begin{document}
\maketitle
\noindent Dorward's Conjecture 4.1 \cite{D} follows by identifying its generating
polynomial with that of a parabolic quotient of the Weyl group of type $B$.
The required unimodality theorem is Stanley's classical Theorem 3.1 \cite{S}.

\begin{theorem}
For every $n\geq3$, the polynomial $\operatorname{LS}_{n,1}(112)$ is unimodal.
More precisely, put $m=n-2$ and $[r]_q=1+q+\cdots+q^{r-1}$. Then
\[
 q^{-1}\operatorname{LS}_{n,1}(112)
 =[m]_q\prod_{i=2}^m(1+q^i)
\]
is symmetric and unimodal of degree $d=(m^2+3m-4)/2$.
\end{theorem}
\begin{proof}
The displayed product formula is Dorward's Theorem 3.3, proved in
Proposition 3.5 \cite{D}; we use it as an attributed input. If $m=1$,
the product is $1$ and the assertion is immediate. Assume $m\geq2$.

Let $W$ be the Weyl group of type $B_m$. Number its simple reflections
$s_1,\ldots,s_m$ along the Dynkin diagram, with $s_m$ at the short-root
end, and set
\[
 J=\{s_1,\ldots,s_{m-2},s_m\}.
\]
Deleting $s_{m-1}$ separates the diagram into types $A_{m-2}$ and $A_1$.
Here $A_0$ is the trivial group when $m=2$.
Consequently the standard parabolic subgroup $W_J$ has type
$A_{m-2}\times A_1$.

For a finite Weyl group $U$, write $F_U(q)=\sum_{w\in U}q^{\ell(w)}$,
where $\ell$ is Coxeter length. The classical length products, with the
empty-product convention, are
\[
 F_W(q)=\prod_{i=1}^m[2i]_q,\qquad
 F_{W_J}(q)=[2]_q\prod_{i=1}^{m-1}[i]_q.
\]
These are the type $B$ and type $A$ formulas in \cite[\S\S3--5]{S}.
Each right coset has a unique element $v$ of minimum length, and every
$w\in W$ has a unique factorization $w=vu$, $u\in W_J$, with
$\ell(w)=\ell(v)+\ell(u)$. Thus the minimum-length representatives
$W^J$ have generating polynomial
\begin{align*}
 F_{W^J}(q)
 &=\frac{\prod_{i=1}^m[2i]_q}
 {[2]_q\prod_{i=1}^{m-1}[i]_q}\\
 &=\frac{[m]_q}{[2]_q}\prod_{i=1}^m(1+q^i)
 =[m]_q\prod_{i=2}^m(1+q^i).
\end{align*}
The cancellation uses $[2i]_q=[i]_q(1+q^i)$ and $[2]_q=1+q$.
Stanley's Theorem 3.1 states that the length coefficients of $W^J$ are
symmetric and unimodal for every parabolic quotient of a finite Weyl
group. This proves the assertion for the displayed product. Its degree is
\[
 (m-1)+\sum_{i=2}^m i=\frac{m^2+3m-4}{2}.
\]
Multiplication by $q$ prepends a zero coefficient, which preserves
unimodality since all coefficients are nonnegative.
\end{proof}

\noindent\textbf{Why the classical theorem applies.}
The quotient above corresponds to the smooth complex projective homogeneous
space $\mathrm{SO}(2m+1,\mathbb C)/P$, where the Levi Weyl group of $P$ is
$W_J$. Its Schubert cells give
$F_{W^J}(q)=\sum_{r=0}^d\dim H^{2r}(G/P,\mathbb C)q^r$.
Hard Lefschetz makes multiplication by an ample class injective from
$H^{2r}$ to $H^{2r+2}$ for $r<d/2$; Poincar\'e duality gives symmetry.
This is Stanley's argument, not a new proof of hard Lefschetz.

\noindent\textbf{Scope and verification.}
For $n<3$ there are no such restricted growth functions, so the polynomial
is zero. We prove the coefficient assertion in Conjecture 4.1. Equality
of rank polynomials does not identify Dorward's poset $M^1(m)$ with
$W^J$, and no Sperner or symmetric-chain result for $M^1(m)$ is claimed.
The accompanying exact verifier independently enumerates restricted growth
functions of lengths $0$ through $9$ and minimum-length signed coset
representatives for $2\leq m\leq10$. Those finite checks support the
identification; the all-parameter proof is the argument above.

\subsection*{Tasks 2--5: proposed follow-ups}
These are proposals, not additional solutions or campaign credits.
\begin{enumerate}
\item \textbf{Task 2.} Decide the Sperner assertion for Dorward's poset
$M^1(m)$, which is not determined by its rank polynomial.
\item \textbf{Task 3.} Construct an explicit statistic-preserving bijection
between the restricted growth functions and the signed coset representatives,
with the one-unit shift in degree included.
\item \textbf{Task 4.} Construct symmetric chain decompositions of $M^1(m)$,
or identify an obstruction to their existence.
\item \textbf{Task 5.} Determine which exact-$k$ occurrence generating
polynomials admit homogeneous-space interpretations or other Lefschetz
operators, starting with $k=2$.
\end{enumerate}

\begin{thebibliography}{9}\raggedright
\bibitem{D} R. Dorward, \emph{Combinatorial statistics on restricted growth
functions containing a pattern exactly $k$ times}, undergraduate thesis
(internally published 2017), arXiv:2003.05073v1 (2020), Theorem 3.3 and
Conjecture 4.1. \url{https://arxiv.org/abs/2003.05073}.
\bibitem{S} R. P. Stanley, \emph{Weyl groups, the hard Lefschetz theorem,
and the Sperner property}, SIAM J. Algebraic Discrete Methods \textbf{1}
(1980), 168--184, Theorem 3.1 and \S\S3--5.
\url{https://math.mit.edu/~rstan/pubs/pubfiles/42.pdf}.
\end{thebibliography}
\end{document}
